% Part: sets-functions-relations
% Chapter: sets
% Section: pairs-and-products

\documentclass[../../../include/open-logic-section]{subfiles}

\begin{document}

\olfileid[tr]{sfr}{set}{pai}
\olsection{İkililer, Demetler ve Kartezyen Çarpımlar}

\begin{explain}
Kaplamsallıktan, bir kümenin elemanları arasında herhangi bir sıra bulunmadığı
sonucu çıkar. Bu yüzden bir sırayı temsil etmek istediğimizde \emph{sıralı
ikilileri} $\tuple{x, y}$ kullanırız. Sırasız bir $\{x, y\}$ ikilisinde sıra
önemli değildir: $\{x, y\} = \{y, x\}$. Sıralı bir ikilide ise önemlidir:
$x \neq y$ ise $\tuple{x, y} \neq \tuple{y, x}$ olur.

Küme kuramında sıralı ikilileri nasıl düşünmeliyiz? En önemlisi, sıralı
ikililerin ancak ve ancak birinci elemanları kendi aralarında, ikinci
elemanları da kendi aralarında aynı olduğunda özdeş olduğu düşüncesini korumak
isteriz; yani
\[
  \tuple{a, b}= \tuple{c, d}\quad\Longleftrightarrow\quad
  a = c\ \text{ ve }\ b=d.
\]
Sıralı ikilileri küme kuramında Wiener--Kuratowski tanımını kullanarak
verebiliriz.
\end{explain}

\begin{defn}[Sıralı ikili]\ollabel{wienerkuratowski}
	$\tuple{a, b} = \{\{a\}, \{a, b\}\}$.
\end{defn}

\begin{prob}
	\olref[sfr][set][pai]{wienerkuratowski} tanımını kullanarak
	$\tuple{a, b}= \tuple{c, d}$ eşitliğinin ancak ve ancak hem $a = c$ hem de
	$b=d$ olduğunda geçerli olduğunu ispatlayın.
\end{prob}

\begin{explain}
Sıralı ikili tanımını belirledikten sonra onu başka kümeleri tanımlamak için
kullanabiliriz. Örneğin bazen ikiden fazla nesnenin sıralı dizilerine de gerek
duyarız: \emph{üçlüler} $\tuple{x, y, z}$, \emph{dörtlüler}
$\tuple{x, y, z, u}$ vb. Üçlüleri, ilk elemanı kendisi de bir sıralı ikili
olan özel sıralı ikililer olarak düşünebiliriz: $\tuple{x, y, z}$,
$\tuple{\tuple{x, y},z}$ demektir. Aynı şey dörtlüler için de geçerlidir:
$\tuple{x,y,z,u}$, $\tuple{\tuple{\tuple{x,y},z},u}$ demektir ve böyle devam
eder. Genel olarak \emph{sıralı $n$-lilerden}
$\tuple{x_1, \dots, x_n}$ söz ederiz.

Sıralı ikililerden veya başka sıralı $n$-lilerden oluşan belirli kümeler
yararlı olacaktır.
\end{explain}

\begin{defn}[Kartezyen çarpım]
$A$ ve $B$ kümeleri verildiğinde bunların \emph{Kartezyen çarpımı}
$A \times B$ şöyle tanımlanır:
\[
  A \times B = \Setabs{\tuple{x, y}}{x \in A \text{ ve } y \in B}.
\]
\end{defn}

\begin{ex}
$A = \{0, 1\}$ ve $B = \{1, a, b\}$ ise çarpımları şöyledir:
\[
A \times B = \{ \tuple{0, 1}, \tuple{0, a}, \tuple{0, b},
    \tuple{1, 1}, \tuple{1, a}, \tuple{1, b} \}.
\]
\end{ex}

\begin{ex}
$A$ bir kümeyse $A$ kümesinin kendisiyle çarpımı $A \times A$, $A^2$ olarak da
yazılır. Bu, $x, y \in A$ koşulunu sağlayan \emph{bütün}
$\tuple{x, y}$ ikililerinin kümesidir. Bütün $\tuple{x, y, z}$ üçlülerinin
kümesi $A^3$'tür ve böyle devam eder. Özyinelemeli bir tanım verebiliriz:
\begin{align*}
  A^1 & = A\\
  A^{k+1} & = A^k \times A
\end{align*}
\end{ex}

\begin{prob}
$\{1, 2, 3\}^3$ kümesinin bütün elemanlarını listeleyin.
\end{prob}

\begin{prop}\ollabel{cardnmprod}
$A$ kümesinin $n$ elemanı, $B$ kümesinin de $m$ elemanı varsa $A \times B$
kümesinin $n\cdot m$ elemanı vardır.
\end{prop}

\begin{proof}
$A$ kümesinin her $x$ elemanı için, $y \in B$ olmak üzere $A \times B$
kümesinde $\tuple{x, y}$ biçiminde $m$ eleman bulunur.
$B_x = \Setabs{\tuple{x, y}}{y \in B}$ olsun. $x_1 \neq x_2$ olduğunda
$\tuple{x_1, y} \neq \tuple{x_2, y}$ de olduğundan
$B_{x_1} \cap B_{x_2} = \emptyset$ olur. Ancak
$A = \{x_1, \dots, x_n\}$ ise
$A \times B = B_{x_1} \cup \dots \cup B_{x_n}$ olur; dolayısıyla bu kümede
$n\cdot m$ eleman bulunur.

Bunu görselleştirmek için $A \times B$ kümesinin elemanlarını bir
ızgarada düzenleyin:
\[
\begin{array}{rcccc}
  B_{x_1} = & \{\tuple{x_1, y_1} & \tuple{x_1, y_2} & \dots & \tuple{x_1, y_m}\}\\
  B_{x_2} = & \{\tuple{x_2, y_1} & \tuple{x_2, y_2} & \dots & \tuple{x_2, y_m}\}\\
  \vdots & & \vdots\\
  B_{x_n} = & \{\tuple{x_n, y_1} & \tuple{x_n, y_2} & \dots & \tuple{x_n, y_m}\}
\end{array}
\]
$x_i$ değerlerinin tümü birbirinden, $y_j$ değerlerinin tümü de birbirinden
farklı olduğundan, bu ızgaradaki ikililerin hiçbiri bir diğeriyle aynı değildir;
toplam $n\cdot m$ ikili vardır.
\end{proof}

\begin{prob}
$k$'ye göre tümevarımla, tüm $k \ge 1$ değerleri için, $A$ kümesinin $n$
elemanı varsa $A^k$ kümesinin $n^k$ elemanı bulunduğunu
gösterin.
\end{prob}

\begin{ex}
$A$ bir kümeyse $A$ üzerindeki bir \emph{sözcük}, $A$ kümesinin
elemanlarından oluşan herhangi bir \emph{sonlu} dizidir. Böyle bir sonlu dizi,
$A$ kümesinin elemanlarından oluşan bir $n$-li olarak düşünülebilir. Örneğin
$A = \{a, b, c\}$ ise ``$bac$'' dizisi $\tuple{b, a, c}$ üçlüsü olarak
düşünülebilir. Sözcükler, yani simge dizileri, bilgisayar biliminde temel bir
öneme sahiptir. Uzlaşım gereği $A$ kümesinin elemanlarını uzunluğu
$1$ olan diziler, $\emptyset$ kümesini de uzunluğu $0$ olan dizi sayarız.
Böylece $A$ üzerindeki \emph{bütün} sözcüklerin kümesi şöyledir:
\[
A^* = \{\emptyset\} \cup A \cup A^2 \cup A^3 \cup \dots
\]
\end{ex}

\end{document}
